% !TeX root = ../main.tex
%------------------------------------------------------------------------------
%  Support and Reinforcement
%  Sources: Hoek, Kaiser & Bawden (1995); Windsor & Thompson (1993);
%  Thomas (2009) on sprayed concrete linings; ACI PRC-506; EN 14487.
%------------------------------------------------------------------------------
\chapter{Support and Reinforcement}
\label{ch:support}

The preceding chapters establish how much load the ground will shed and how
much convergence it will undergo. This one deals with the hardware that
receives that load: bolts and cables acting within the rock mass, sprayed
concrete acting on its surface, and mesh and plates holding the loose material
between them.

Three ideas organise the chapter. The first is that reinforcement and support
are mechanically distinct and are analysed differently. The second is that
what a support element contributes is set by its stiffness, its capacity and
the convergence that had already occurred when it was installed --- the
argument developed in Chapter~\ref{ch:ccm} and applied here to real hardware.
The stiffness and capacity expressions for each element are collected in
Section~\ref{sec:grc-scc-appendix}, representative properties in
Appendix~\ref{app:constants}, and the sequence of the interaction
calculation in Appendix~\ref{app:calculations}.
The third is that installation quality determines as-built performance more
often than design does. A correctly designed bolt that is not tensioned, or a
shotcrete layer 40\,mm thinner than specified over half its area, has failed
in a way that an analysis run on nominal dimensions does not capture.

\section{Reinforcement principles and hardware}
\label{sec:reinforcement-principles}

\subsection{Reinforcement, support and retention}
\label{sec:reinforcement-vs-support}

Three functions are routinely conflated in specifications and must be
separated in analysis.

\emph{Reinforcement} acts within the rock mass, improving its own capacity. A
grouted bolt crossing a discontinuity adds confinement across it, raises the
frictional resistance available on it, and resists shear across it. The rock
mass remains the load-carrying structure; the reinforcement makes it better at
the job.

\emph{Support} acts on the excavation boundary, applying a pressure that
resists convergence. Shotcrete shells, steel sets and concrete linings are
support: they are structures in their own right, carrying load that the ground
has shed onto them.

\emph{Retention} holds loose material in place between other elements. Mesh
and straps do almost nothing for global stability but prevent the small
falls of ground that account for most underground
injuries~\citep{hoek1995support}.

The distinction determines the analysis. Reinforcement is assessed through its
effect on the rock mass --- higher effective confinement, a raised ground
reaction curve, a stabilised block. Support is assessed through its own
structural capacity against the pressure it attracts. A design that treats
grouted bolts as though they applied a boundary pressure has used the wrong
model, and one that omits retention has designed a system that will still
injure people.

\subsection{Bolt types and load transfer}
\label{sec:bolt-types}

Reinforcement systems are usefully classified by how load transfers from the
rock to the element~\citep{windsor1993classification}. \emph{Continuously mechanically coupled} systems --- fully
grouted rebar, resin-anchored bar --- transfer load along their whole length
through a stiff grout or resin annulus. \emph{Continuously frictionally
coupled} systems --- split sets, expandable friction bolts --- transfer through
friction at the borehole wall along their length. \emph{Discretely mechanically
coupled} systems --- mechanical expansion-shell anchors, and cables anchored
only at their ends --- transfer at discrete points.

The classification predicts behaviour. Continuously coupled elements develop
load locally wherever the rock moves, which makes them stiff, effective across
a discontinuity that opens anywhere along the bolt, and tolerant of losing any
one section. Discretely coupled elements develop load only through relative
displacement between their anchor points, which makes them softer and reliant
on the anchor remaining sound. Frictionally coupled elements are the softest
and the most forgiving of large deformation, which is why they persist in
squeezing and burst-prone ground despite modest capacity. Representative
capacities for each family are in Table~\ref{tab:support-props} of
Appendix~\ref{app:constants}.

Load transfer along a grouted bolt is not uniform. Shear stress concentrates
near the point of relative movement and decays along the bond, so only a
limited embedment length is mobilised at any time. Past a critical shear
stress the bond decouples progressively, and the load redistributes further
along the bar. Two design consequences follow: additional embedment does not
raise the load at which the bond first mobilises, but it provides the reserve
into which load redistributes as the near-field bond decouples; and grout
quality along the whole
bond governs performance in a way a pull test at one location may not reveal.

Self-drilling anchors deserve mention because they solve a construction
problem rather than a mechanical one. Where a borehole will not stay open ---
in ravelling ground, in heavily fractured rock, in soil --- a hollow bar
drilled and grouted in one operation is the only practical option, at some
cost in steel efficiency.

\subsection{Design for the governing mechanism}
\label{sec:design-mechanism}

Bolting design differs fundamentally according to what is being resisted, and
the commonest error is to apply the pattern appropriate to one mechanism where
the other governs.

For \emph{structurally controlled} failure the calculation is discrete. The
block is identified by the methods of
Chapter~\ref{ch:structural-analysis}, its weight and sliding mode established (Section~\ref{sec:block-and-wedge-construction}),
and the bolt force required to reach a target factor of safety follows from
limit equilibrium. Bolt orientation matters: the optimum is generally inclined
to the sliding direction rather than normal to the excavation surface, because
an inclined bolt contributes both a direct restraining component and an
increase in normal stress on the sliding plane. Bolts must be long enough to
anchor beyond the block, which for a large crown wedge can exceed the length
that standard equipment can install.

For \emph{stress-induced} failure the calculation is continuum. Systematic
bolting supplies a modest confining pressure at the boundary, and modest is
enough: rock strength is strongly confinement-dependent near a free surface,
so a few tens to about a hundred kilopascals of confinement produces a
disproportionate increase in strength. The bolting must extend beyond the computed yield zone
(Chapter~\ref{ch:ground-behaviour}) so that it is anchored in rock that has
not failed, and its effect appears as a raised ground reaction curve rather
than as a resisted block weight.

Pattern bolting reconciles the two in practice. Length is set from the larger
of the yield zone depth and the critical block dimension; spacing from the
block size that must be retained and from the confinement required. Empirical
charts give both directly from classification (Chapter~\ref{ch:empirical-design}) and
remain the sensible starting point, with the analytical checks used to confirm
that the chart's implicit mechanism is the one actually operating.

\section{Reinforcement design}
\label{sec:reinforcement-design}

\subsection{Shear, cables and dynamic reinforcement}
\label{sec:shear-cables-dynamic}

A bolt crossing a discontinuity that shears is loaded in combined tension and
shear, not tension alone. As the joint displaces, the bar bends near the
joint, develops axial tension through the geometry of that deformation, and
contributes dowel action --- direct shear resistance across the plane. The
combination is complex, the bar frequently fails at the joint where bending
and tension coincide, and the practical consequence is that a bolt oriented
perpendicular to a shearing joint is poorly used. An inclination of roughly
$20$ to $30$ degrees from the normal, in the direction that puts the bar into
tension as the joint moves, is substantially more efficient.

Cables extend reinforcement to lengths bolts cannot reach. Plain strand,
bulbed strand and multi-strand anchors span the range from a few metres to
tens of metres, and pre-tensioning is used where the objective is to prevent
displacement rather than to resist it after it occurs. Their role is in large
spans --- cavern roofs and high sidewalls (Chapter~\ref{ch:openings}) --- where
the blocks that can form are larger than any bolt can reach past. Because
cables are usually discretely coupled, they are softer than grouted bolts, and
a system combining both will load the bolts first.

Burst-prone ground requires a different criterion entirely. The demand is
energy rather than force: rock ejected at some velocity carries kinetic energy
that the support must absorb, and a stiff high-capacity bolt that cannot
deform absorbs almost none of it. Energy-absorbing bolts yield at a controlled
load over a substantial displacement, and their capacity is quoted in
kilojoules rather than kilonewtons.

System compatibility is the governing constraint here, and it is frequently
missed: every element has to deform through the same displacement, and in
dynamic ground the retention system, not the bolt, is usually the weak link.
The energy demand and the matching of elements are set out in
Chapter~\ref{ch:high-stress}.

\subsection{Surface retention and durability}
\label{sec:retention-durability}

Mesh is the least analysed and most frequently damaged element of a support
system. Welded mesh is stiff and predictable but fails brittly at its welds;
chain-link is far more ductile and can sustain large deformation without
losing retention, which is what squeezing and burst-prone ground demand. What matters
mechanically is the connection: mesh performs only as well as its attachment
to the bolts, and the plate and its fixing are as much part of the system as
the mesh sheet.

Plate stiffness governs how load is collected from the surface into the bolt.
A thin dished plate on a rough surface bears at a few points and deforms
locally rather than transferring the load. This is a detailing question with
mechanical consequences, and it is the kind of detail that is value-engineered
out of a specification without anybody analysing what was lost.

Durability separates temporary from permanent reinforcement. A bolt required
for the construction period only can be plain steel. A bolt forming part of
the permanent works must survive a century in ground whose aggressivity is
established from the groundwater chemistry of
Chapter~\ref{ch:hydrogeology}. Galvanising, epoxy coating, full encapsulation
and double corrosion protection form a hierarchy of increasing cost and
reliability. The decision needs to be explicit, because the common default ---
installing temporary-grade reinforcement and later declaring it permanent
because it is still there --- is not a design.

\subsection{Testing and quality control}
\label{sec:bolt-testing}

Bolt testing answers a narrower question than is often assumed. A pull test
loads the bar against the rock surface and demonstrates that the bond over
some length can carry the applied load. It does not demonstrate that the grout
column is complete, that the bar is centred, or that the anchorage beyond the
tested length is sound.

A useful regime therefore combines proof loading of a defined proportion of
installed bolts, destructive testing of a small number to establish the
capacity actually achievable in that ground, grout quality verification
through mix records and cube testing, and installation records for every bolt
--- location, orientation, length, grout batch, date. Overcoring a small
sample to inspect the grout column is the only direct check of encapsulation,
and it is worth doing where the reinforcement is permanent.

Acceptance criteria must be set against the design assumption rather than
against the bar's yield capacity alone. Bond governs wherever embedment,
grout quality or the condition of the borehole limits load transfer, so the
criterion is set against the lower of the design bond capacity and the bar
capacity, with the governing mode identified for the ground in question.

\section{Sprayed concrete}
\label{sec:sprayed-concrete}

\subsection{Role and process}
\label{sec:shotcrete-role}

Sprayed concrete began as a surface sealant, protecting weathering-sensitive
rock and preventing small falls. It has become the primary structural lining
for most conventionally excavated tunnels and, increasingly, a permanent one.
That progression matters because the analysis, testing and quality regime
appropriate to a sacrificial coating is inadequate for a structure carrying
permanent load.

The wet-mix process, in which fully batched concrete is pumped to the nozzle
and accelerator added there, has displaced dry-mix almost entirely on major
works. It gives better control of water content, far less rebound and dust,
and much higher output. Dry-mix retains a role for small quantities,
intermittent work and repairs, where the ability to stop and start without
wasting a batch outweighs its disadvantages.

Rebound is real material loss and a real cost --- and it is selective, since
coarse aggregate rebounds preferentially and the material that stays on the
wall is richer in fines and cement than the mix that left the pump. Rebound
should be measured, not assumed, and mix design should account for the
composition that actually adheres.

\subsection{Mix design and reinforcement}
\label{sec:shotcrete-mix}

Sprayed concrete must satisfy two conflicting requirements: it must be pumpable
and sprayable, which favours a fine, cohesive, workable mix, and it must
develop strength within hours, which conflicts with the water content that
pumpability demands. The mix is a negotiation between them, mediated by
plasticisers, hydration control admixtures that extend working life, and
accelerators added at the nozzle.

Accelerator choice has consequences beyond setting time. Alkali-free
accelerators have displaced alkaline products for health reasons --- the older
products caused severe burns --- and they also reduce the long-term strength
loss that heavy accelerator dosage causes~\citep{thomas2009scl}. That loss is
typically 15 to 30 per cent relative to the unaccelerated mix, and it must be
allowed for in the design strength rather than discovered in core testing.

Fibre reinforcement has largely replaced mesh in sprayed linings, and the
reason is constructability as much as structure: fixing mesh to an irregular
profile is slow, hazardous and produces shadowing behind the bars where the
concrete does not compact. Steel fibres give higher residual strength and
stiffness; macro-synthetic fibres are lighter, non-corroding and easier to
handle, at lower residual capacity. What is specified is not a dosage but a
performance --- residual flexural strength at defined crack openings, or
energy absorption from a panel test --- because fibre performance depends on
dosage, orientation and distribution together, and only a test on sprayed
material measures the combination.

\subsection{Early-age and structural behaviour}
\label{sec:shotcrete-early-age}

Early-age strength, not 28-day strength, governs tunnelling. The excavation
cycle cannot advance until the lining can carry load and cannot be worked
beneath until it is safe, so the strength development curve over the first few
hours determines the cycle time and therefore the programme. Indicative
strength and stiffness development over the first day is in
Appendix~\ref{app:constants}.

The early strength classes distinguish rapid-setting mixes for overhead work
in poor ground from slower mixes acceptable in better conditions, and
compliance is verified in the tunnel by penetrometer at very early ages and
stud driving thereafter. These are crude tests, and they are the ones the
decision to re-enter actually rests on.

Structurally, a sprayed lining is a thin shell in which thrust dominates and
moment is secondary. Capacity is therefore assessed on a thrust--moment
interaction diagram, the locus of the pairs of axial thrust and bending
moment the section can carry, so that a thrust $N$ acting at an eccentricity
$e$ plots at $M = Ne$ and must fall inside the envelope
(Chapter~\ref{ch:linings}). It is not assessed against a bending capacity
alone, and the
governing case is often a modest thrust with an eccentricity caused by
profile irregularity rather than the design thrust at the design thickness.

Two time-dependent effects, whose histories run in opposite directions, are
frequently neglected. Stiffness rises as the concrete cures, so a lining installed
against a converging ground attracts load progressively rather than all at
once. Creep relaxes that load, and at early ages the creep of young shotcrete
is very large --- enough that the load computed from an elastic shell with
mature stiffness may exceed the actual load by a factor of two. Ignoring both
overestimates the load on the lining, which is safe for its capacity check,
while underestimating the convergence that the clearance envelope and any
movement limit on adjacent assets must accommodate --- which is why load
estimated from measured convergence is more reliable than load computed from
an assumed stiffness.

\section{Sprayed lining design and control}
\label{sec:scl-design}

\subsection{Lining design}
\label{sec:scl-design-approach}

The first decision is whether the sprayed lining is temporary, permanent, or
part of a composite. A \emph{sacrificial} primary lining carries construction
loads and is assumed to contribute nothing thereafter, with a secondary lining
designed for the full long-term load. A \emph{composite} lining assumes the
two act together, which requires the interface to transfer shear and the
primary lining's long-term durability to be demonstrated. A \emph{single-shell}
lining is the permanent structure in its own right.

Each is defensible; the differences are in verification burden. Single-shell
and composite designs are more economical and lower in carbon, and they
require thickness control, durability evidence and quality records that a
sacrificial lining does not. Choosing the economical option without accepting
the verification obligation is the failure mode to avoid.

Squeezing ground requires the lining to be allowed to deform. A stiff shell
installed early in convergent ground attracts a load it cannot carry, and the
result is systematic cracking and reprofiling. Deformation slots --- gaps in
the shell fitted with yielding elements that shorten at controlled load ---
let the lining converge while continuing to provide support pressure, so
equilibrium is reached on the ground reaction curve at a pressure the lining
can sustain. This is the clearest practical demonstration that support design
is about equilibrium at an acceptable convergence, not about minimising
convergence (Chapter~\ref{ch:ccm}).

Reinforced ribs of sprayed concrete offer an intermediate stiffness between
plain shotcrete and steel sets: reinforcing bars laid circumferentially and
encased in a thickened sprayed rib, giving substantial bending capacity
without the fit-up problems of steel arches.

\subsection{Steel elements and substrate}
\label{sec:steel-and-substrate}

Lattice girders and steel sets serve three purposes that are worth separating
because they have different design implications. They control profile, giving
the crew a physical template for thickness. They provide immediate support in
the interval before the shotcrete gains strength. And they contribute bending
capacity to the completed lining.

The third is the contested one. A lattice girder fully encased in sprayed
concrete of comparable modulus contributes far less than its section
properties suggest, and its principal value is the first two functions. Heavy
steel arches genuinely contribute capacity but are difficult to encase without
voids behind the flanges, and those voids are both a durability defect and a
loss of the composite action assumed.

Substrate preparation determines whether the lining acts as a shell at all.
Bond to the rock and between successive layers requires a clean surface free
of loose material, dust and standing water. Where bond fails, layers act
independently and the effective thickness is that of one layer rather than
their sum --- a reduction in bending capacity by a factor of several. Cold
joints between layers sprayed hours or days apart are the most common
location.

Water inflow through the substrate is the most difficult condition. Water
prevents bond, washes out cement paste locally and creates a permanent path.
Local drainage --- pipes and fleece strips leading water to a collection point
--- must be installed before spraying, and accepting a wet spot and spraying
over it produces a defect that reappears as a leak.

\subsection{Durability and quality control}
\label{sec:scl-durability-qc}

A sprayed permanent lining must satisfy the same durability requirements as
cast concrete, against the exposure established in
Chapter~\ref{ch:hydrogeology}, and it must do so with a material placed by
projection onto an irregular surface. Cover to reinforcement, permeability and
crack width are the controlling parameters, and all three are more variable in
sprayed than in cast concrete.

Fire performance requires explicit attention. Dense low-permeability concrete
spalls explosively under rapid heating as pore water vaporises faster than it
can escape, and the mixes used for durability are precisely the ones most
vulnerable. Polypropylene micro-fibres, which melt and create escape paths,
are the standard mitigation and must be specified for any lining with a fire
requirement (Chapter~\ref{ch:linings}).

Quality control is where sprayed concrete design most often fails. Thickness
is the variable that matters and the hardest to control: it is placed by eye
against an irregular profile, and under-thickness is systematic rather than
random because a nozzleman working overhead against gravity naturally applies
less. Depth pins, laser scanning of the sprayed surface against the design
profile, and coring for verification are all used, and scanning is now the
practical standard because it measures the whole surface rather than a sample.

Strength verification requires cores from sprayed panels or from the lining
itself, not cubes cast from the pump. Material that has been through the
nozzle differs from the batch that entered it, and testing the batch measures
something that was never installed.

\chapterreferences
